\section{Conclusions and operational implications}
Accepted adaptation is an expansion of an optimized contractual opportunity
set, not permission to disregard a commitment. The tree results identify that
expansion through continuation prices, switching capacities, and constructive
deviations. Benchmark-relative transfer coordinates make the connection exact
for positive scalar payments: their lower bounds retain the participation
slack of a nonconstant optimized contract. Cap matching recovers the original
nonnegative flows; it must not be imposed on a benchmark that has already
earned unused continuation slack. The quadratic formula and exact example show
how this distinction changes the economically relevant expansion value.

Implementation introduces a separate loss. Resource-value gradients,
direct-price approximations, and classical continuation can propose decisions
on the same accepted geometry. Under the stated strong convexity and smooth
coupling, the global bridge controls price-induced regret across faces and
switching changes. General convex objectives retain their separate certificates;
adding regularization does not remove the original-objective bias charge.
Component repair identifies fixed switching and participation floors, whereas
complete dual repair attains the true regret floor without strict feasibility.
A better upper certificate does not itself improve the underlying decision.

The empirical results preserve this separation. Optimized restricted
comparators measure accepted expansion; exact audits measure implementation
loss; repeated-query accounting measures whether approximation saves work.
Paid optimizer labels and fallback are part of that account, not free inputs.
The robust certification lemma protects the model-specific comparator inside a
declared coefficient set, and direct interior-model brackets quantify its
conservatism. A designed coefficient cube is not a calibrated uncertainty set,
and a frozen-policy guarantee is not a theorem about arbitrary retraining.
The retained neural, direct-price, and classical comparisons support an
accepted-control framework with interchangeable implementations, rather than a
claim that one predictor must dominate every decision regime.
